\subsection{Conditioning}

Throughout this subsection, $\mathbb{F}\in\{\mathbb{R},\mathbb{C}\}$, and all matrix dimensions are positive integers.

\begin{definition} \label{def:spectral-condition-number} \hlt{Spectral Condition Number}\\
If $A\in\operatorname{Mat}(n,n,\mathbb{F})$ is invertible, its \hlt{spectral condition number} is $\kappa_2(A)\defeq\norm{A}_2\norm{A^{-1}}_2$. If $A$ is singular, $\kappa_2(A)$ is undefined.
\end{definition}

\begin{proposition} \label{prop:spectral-condition-number-properties} \hlt{Properties of the Spectral Condition Number}\\
Let $A,B\in\operatorname{Mat}(n,n,\mathbb{F})$ be invertible and let $c\in\mathbb{F}\setminus\{0\}$. Let $U,V\in\operatorname{Mat}(n,n,\mathbb{F})$ be unitary.
\begin{enumerate}[label=(\roman*)]
\setlength{\itemsep}{0pt}
\item $\kappa_2(AB)\leq\kappa_2(A)\kappa_2(B)$;
\item $\kappa_2(A)=\frac{\sigma_{\max}(A)}{\sigma_{\min}(A)}\geq1$, with equality if and only if $A$ is a nonzero scalar multiple of a unitary matrix;
\item $\kappa_2(cA)=\kappa_2(A)$;
\item $\kappa_2(A)=\kappa_2(A^*)=\kappa_2(A^{-1})$; over $\mathbb{C}$, these also equal $\kappa_2(\overline{A})$ and $\kappa_2(A^t)$;
\item $\kappa_2(A^*A)=\kappa_2(A)^2$;
\item $\kappa_2(UAV)=\kappa_2(A)$.
\end{enumerate}
\begin{proof}
\begin{enumerate}[label=(\roman*)]
\setlength{\itemsep}{0pt}
\item[$(i)$] By Proposition \ref{prop:spectral-norm-basic-properties}(v), $\kappa_2(AB)=\norm{AB}_2\norm{B^{-1}A^{-1}}_2\leq\kappa_2(A)\kappa_2(B)$.
\item[$(ii)$] ($\Rightarrow$) Fix SVD $A=X\Sigma Y^*$ by Theorem \ref{thm:matrix-singular-value-decomposition}. Since $A$ invertible, $\sigma_n(A)>0$ by Corollary \ref{cor:square-invertibility-singular-value-dictionary}(i)$\Rightarrow$(iii), $A^{-1}=Y\Sigma^{-1}X^*$.\\
Thus $\kappa_2(A)=\frac{\sigma_1(A)}{\sigma_n(A)}\geq1$ by Theorem \ref{thm:spectral-norm-largest-singular-value}(i). For $s\defeq\sigma_1(A)>0$, equality forces $\Sigma=sI_n$, so $A=sXY^*$, a nonzero scalar multiple of a unitary matrix by Theorem \ref{thm:unitary-matrix-equivalent-conditions}(iv) and Proposition \ref{prop:algebra-of-unitary-matrices}(i).
\item[$(ii)$] ($\Leftarrow$) If $A=dQ$ with $d\neq0$ and $Q$ unitary, then $Q^*Q=I_n$, so the singular values of $Q$ are all $1$.\\
Those of $A$ are all $\abs{d}$ by Proposition \ref{prop:basic-singular-value-identities}(iv); hence the quotient in part (ii) is $1$.
\item[$(iii)$] Since $(cA)^{-1}=c^{-1}A^{-1}$, $\kappa_2(cA)=\abs{c}\norm{A}_2\abs{c}^{-1}\norm{A^{-1}}_2=\kappa_2(A)$ by Proposition \ref{prop:spectral-norm-basic-properties}(iii).
\item[$(iv)$] The singular values are unchanged under adjoint and, over $\mathbb{C}$, under transpose and conjugation by Proposition \ref{prop:basic-singular-value-identities}(iii).\\
Inversion reverses and reciprocates them by part (ii), so every stated condition number equals $\kappa_2(A)$ by part (ii).
\item[$(v)$] Matrix $A^*A$ positive definite by Corollary \ref{cor:square-invertibility-singular-value-dictionary}(i)$\Leftrightarrow$(iv). Its eigenvalues are $\sigma_1(A)^2,\ldots,\sigma_n(A)^2$ by Proposition \ref{prop:basic-singular-value-identities}(i).\\
These are also its singular values by Theorems \ref{thm:real-spectral-theorem} and \ref{thm:complex-spectral-theorem}, as appropriate, and Proposition \ref{prop:unitary-invariance-of-singular-values}(i).\\
Hence $\kappa_2(A^*A)=\frac{\sigma_1(A)^2}{\sigma_n(A)^2}=\kappa_2(A)^2$ by part (ii).
\item[$(vi)$] The equality $\kappa_2(UAV)=\kappa_2(A)$ holds by Proposition \ref{prop:unitary-invariance-of-singular-values}(i) and part (ii).
\end{enumerate}
\hlt{Structural Commentary}: Norm laws and singular-value transformations transfer directly to $\kappa_2$.\\
\hlt{Key Insight}: Only the ratio $\frac{\sigma_1(A)}{\sigma_n(A)}$ matters.
\end{proof}
\end{proposition}

\begin{corollary} \label{cor:positive-definite-condition-number-eigenvalues} \hlt{Positive Definite Condition Number by Eigenvalues}\\
Let $A\in\operatorname{Mat}(n,n,\mathbb{F})$ satisfy $A\succ0$ and have eigenvalues $0<\lambda_1\leq\cdots\leq\lambda_n$.
\begin{enumerate}[label=(\roman*)]
\setlength{\itemsep}{0pt}
\item $\norm{A}_2=\lambda_n$;
\item $\norm{A^{-1}}_2=\lambda_1^{-1}$;
\item $\kappa_2(A)=\frac{\lambda_n}{\lambda_1}$.
\end{enumerate}
\begin{proof}
\begin{enumerate}[label=(\roman*)]
\setlength{\itemsep}{0pt}
\item[$(i)$] The matrix $A$ is self-adjoint with positive eigenvalues by Theorem \ref{thm:positive-definite-matrix-equivalences}(i)$\Rightarrow$(ii).\\
Singular values are $\lambda_n,\ldots,\lambda_1$ by Theorems \ref{thm:real-spectral-theorem}, \ref{thm:complex-spectral-theorem}, as appropriate, Proposition \ref{prop:unitary-invariance-of-singular-values}(i); $\norm{A}_2=\lambda_n$ by Theorem \ref{thm:spectral-norm-largest-singular-value}(i).
\item[$(ii)$] The same diagonalization has diagonal entries $\lambda_1^{-1},\ldots,\lambda_n^{-1}$ for $A^{-1}$, so these are its singular values by Proposition \ref{prop:unitary-invariance-of-singular-values}(i).\\
Thus $\norm{A^{-1}}_2=\lambda_1^{-1}$ by Theorem \ref{thm:spectral-norm-largest-singular-value}(i).
\item[$(iii)$] The identity $\kappa_2(A)=\frac{\lambda_n}{\lambda_1}$ follows by parts (i)--(ii).
\end{enumerate}
\hlt{Structural Commentary}: Diagonalization assigns the two operator norms to opposite endpoint eigenvalues.\\
\hlt{Key Insight}: Positive definiteness removes absolute values and makes $\lambda_1$ nonzero.
\end{proof}
\end{corollary}

\begin{perspective} \label{per:determinant-is-not-conditioning} \hlt{Determinant is Not Conditioning}\\
Let $A\in\operatorname{Mat}(n,n,\mathbb{F})$ invertible, $c\in\mathbb{F}\setminus\{0\}$. Then $\kappa_2(cA)=\kappa_2(A)$ by Proposition \ref{prop:spectral-condition-number-properties}(iii); $\det(cA)=c^n\det(A)$ by Theorem \ref{thm:determinant-basic-laws}(ii).\\
Thus determinant records volume scaling, not sensitivity under uniform rescaling.
\end{perspective}

\begin{theorem} \label{thm:linear-system-relative-error-condition-number} \hlt{Condition Number and Linear-System Error}\\
Let $A\in\operatorname{Mat}(n,n,\mathbb{F})$ be invertible, $\mathbf{x},\Delta\mathbf{x},\Delta\mathbf{y}\in\mathbb{F}^n$, $\mathbf{y}\in\mathbb{F}^n\setminus\{\mathbf{0}\}$. If $A\mathbf{x}=\mathbf{y}$ and $A\Delta\mathbf{x}=\Delta\mathbf{y}$, then $\frac{\norm{\Delta\mathbf{x}}_2}{\norm{\mathbf{x}}_2}\leq\kappa_2(A)\frac{\norm{\Delta\mathbf{y}}_2}{\norm{\mathbf{y}}_2}$.
\begin{proof}
Since $A$ is invertible and $\mathbf{y}\neq\mathbf{0}$, $\mathbf{x}\neq\mathbf{0}$.\\
By defining norm bound, $\norm{\mathbf{y}}_2\leq\norm{A}_2\norm{\mathbf{x}}_2$, $\norm{\Delta\mathbf{x}}_2=\norm{A^{-1}\Delta\mathbf{y}}_2\leq\norm{A^{-1}}_2\norm{\Delta\mathbf{y}}_2$. Claim follows by division and multiplication.\\
\hlt{Structural Commentary}: The equations bound $\norm{\mathbf{x}}_2$ through $A$ and $\norm{\Delta\mathbf{x}}_2$ through $A^{-1}$.\\
\hlt{Key Insight}: The opposing operator norms multiply to $\kappa_2(A)$.
\end{proof}
\end{theorem}

\begin{theorem} \label{thm:matrix-perturbation-relative-error-condition-number} \hlt{Condition Number and Matrix Perturbation}\\
Let $A\in\operatorname{Mat}(n,n,\mathbb{F})$ be invertible, let $\Delta A\in\operatorname{Mat}(n,n,\mathbb{F})$, and let $\mathbf{x},\Delta\mathbf{x},\mathbf{y}\in\mathbb{F}^n$.\\
Suppose $A\mathbf{x}=\mathbf{y}$, $(A+\Delta A)(\mathbf{x}+\Delta\mathbf{x})=\mathbf{y}$, and $\mathbf{x}+\Delta\mathbf{x}\neq\mathbf{0}$. Then $\frac{\norm{\Delta\mathbf{x}}_2}{\norm{\mathbf{x}+\Delta\mathbf{x}}_2}\leq\kappa_2(A)\frac{\norm{\Delta A}_2}{\norm{A}_2}$.
\begin{proof}
Set $\mathbf{z}\defeq\mathbf{x}+\Delta\mathbf{x}\neq\mathbf{0}$. The system equations imply $\Delta\mathbf{x}=-A^{-1}\Delta A\mathbf{z}$.\\
Hence $\norm{\Delta\mathbf{x}}_2\leq\norm{A^{-1}}_2\norm{\Delta A}_2\norm{\mathbf{z}}_2$ by the defining norm bound and Proposition \ref{prop:spectral-norm-basic-properties}(v).\\
Since $\norm{A}_2>0$ by Proposition \ref{prop:spectral-norm-basic-properties}(ii), the result follows by division by $\norm{\mathbf{z}}_2$ and the defining formula for $\kappa_2(A)$.\\
\hlt{Structural Commentary}: Subtraction isolates $\Delta\mathbf{x}$ as $-A^{-1}\Delta A$ acting on the perturbed solution.\\
\hlt{Key Insight}: Matrix perturbation enters through the product $A^{-1}\Delta A$.
\end{proof}
\end{theorem}

Figure \ref{fig:condition-number-error-bound} depicts the ellipse determined by the extreme singular directions.

\begin{illustrationfigure}[H]
  \centering
  \begin{tikzpicture}[
  >=Stealth,
  every node/.style={font=\small},
  axis/.style={->, gray!70},
  vec/.style={->, line width=1pt},
  outline/.style={draw=gray!70, line width=0.8pt},
  map/.style={->, line width=0.9pt, blue!70!black}
]
  \begin{scope}[shift={(0,0)}]
    \draw[axis] (-1.4,0) -- (1.55,0);
    \draw[axis] (0,-1.15) -- (0,1.35);
    \draw[outline] (0,0) circle (1);
    \draw[vec, teal!75!black] (0,0) -- (0.95,0);
    \draw[vec, red!70!black] (0,0) -- (0,0.95);
  \end{scope}

  \begin{scope}[shift={(4.8,0)}, rotate=18]
    \draw[axis] (-2.05,0) -- (2.2,0);
    \draw[axis] (0,-0.85) -- (0,1.0);
    \draw[outline] (0,0) ellipse (1.8 and 0.45);
    \draw[vec, teal!75!black] (0,0) -- (1.8,0) node[below right] {$\sigma_{\max}$};
    \draw[vec, red!70!black] (0,0) -- (0,0.45) node[above left] {$\sigma_{\min}$};
  \end{scope}

  \draw[map] (1.65,0.25) -- (3.0,0.25) node[midway, above] {$A$};
\end{tikzpicture}

  \caption{The unit circle maps to an ellipse with extreme-semiaxis ratio $\kappa_2(A)=\frac{\sigma_{\max}(A)}{\sigma_{\min}(A)}$.}
  \label{fig:condition-number-error-bound}
\end{illustrationfigure}
